%-----------------------------------------------------------------------------%
\section{Aplikasi \textit{Knowledge Graph} Peraturan Perundang-undangan}
\label{sec:latex}
%-----------------------------------------------------------------------------%
Dibuat chatbot untuk mensimulasikan pengaplikasian KG.
Chatbot akan memproses input yand diberikan oleh user, kemudian meng-generate SPARQL query.
Chatbot akan menampilkan hasil dari query tersebut.

%-----------------------------------------------------------------------------%
\section{Evaluasi Hasil Konversi}
\label{sec:latex}
%-----------------------------------------------------------------------------%
Evaluasi akan dilakukan terhadap semua dokumen yang terdapat pada peraturan.go.id.
Selanjutnya dokumen-dokumen ini akan disebut dokument tes.

%-----------------------------------------------------------------------------%
\section{Evaluasi Manual}
\label{sec:latex}
%-----------------------------------------------------------------------------%
Evaluasi kualitatif dilakukan terhadap sebagian dokumen yang didapatkan dari hasil sampling dari dokumen tes.
Seorang penilai akan melihat markdown dari setiap dokumen hasil \textit{sampling}, masing-masing dokumen akan diberikan skor 0 sampai 100.
Kualitas konversi dapat dilihat dari rata-rata skor tersebut.

%-----------------------------------------------------------------------------%
\section{Evaluasi Otomatis}
\label{sec:latex}
%-----------------------------------------------------------------------------%
Evaluasi dilakukan dengan terlebih dahulu mendefinisikan metadata yang pasti terdapat pada sebuah dokumen (misal, tempat dokumen disahkan), kemudian menghitung seberapa banyak metadata yang berhasil diekstrak.
Metadata hasil ekstraksi yang akan dihitung tidak perlu sesuai dengan yang terdapat di dalam dokumen.
